\documentclass[12pt,a4paper]{article}
\usepackage[utf8]{inputenc}
\usepackage[russian,english]{babel}
\usepackage{amsmath,amssymb,amsthm}
\usepackage{geometry}
\usepackage{hyperref}
\geometry{margin=2.5cm}

\title{
    \textbf{Meta-Stable Architectures}\\
    \Large Volume III (v2): Homeodynamic Memory Regulation and Stochastic Robustness
}
\author{Symbiosis: Katia \& Jippi}
\date{2026}

\newtheorem{theorem}{Theorem}

\begin{document}
\maketitle

\begin{abstract}
Volume III (v2) closes the trilogy with memory-driven homeodynamic regulation under stochastic perturbations. The suppression buffer is constrained to $[0,1]$, and mean-square boundedness is stated with explicit theorem-level assumptions.
\end{abstract}

\selectlanguage{russian}
\begin{abstract}
Том III (v2) завершает трилогию моделью гомеодинамической регуляции памяти при стохастических возмущениях. Буфер подавления ограничен интервалом $[0,1]$, а среднеквадратическая ограниченность сформулирована как строгая теорема с явными предпосылками.
\end{abstract}
\selectlanguage{english}

\section*{Version Notes (v2)}
\begin{itemize}
    \item Added explicit mean-square boundedness theorem assumptions.
    \item Constrained suppression buffer dynamics to an invariant interval $S\in[0,1]$.
    \item Clarified parameter positivity and drift-dissipativity conditions.
    \item Aligned crisis-memory notation with Volumes I--II.
\end{itemize}

\section{Deterministic Multiscale Core}
State vector: $X=(x,y,H,K,\mu,S)\in\mathbb{R}^5\times[0,1]$.
\begin{align}
\dot x &= -k x-\lambda y, \\
\dot y &= \alpha x+\mu y-\gamma y^3, \quad \gamma>0, \\
\dot H &= \sigma_H y^2-(\delta_H+\eta S)H, \\
\dot K &= \beta H-\delta_K K+\xi SH, \\
\dot\mu &= \rho_1 H-\rho_2 K-\rho_3\mu, \\
\dot S &= \epsilon_s(H-H_{\mathrm{crit}})S(1-S).
\end{align}
All coefficients are assumed nonnegative, with
$k,\lambda,\gamma,\delta_H,\delta_K,\rho_3,\epsilon_s>0$.

\section{Stochastic Extension}
Environmental perturbation is added in It\^o form:
\begin{equation}
dy=(\alpha x+\mu y-\gamma y^3)dt+\sigma dW_t.
\end{equation}

\section{Mean-Square Boundedness}
Let
\begin{equation}
\mathcal V(X)=a_1x^2+a_2y^2+a_3H^2+a_4K^2+a_5\mu^2+a_6S^2,
\quad a_i>0.
\end{equation}

\begin{theorem}
Assume local Lipschitz coefficients and existence of constants $c_1,c_0,c_\sigma>0$ such that the It\^o generator satisfies
\begin{equation}
\mathcal L\mathcal V(X)\le -c_1\mathcal V(X)+c_0+c_\sigma\sigma^2
\end{equation}
for all $X$. Then, for every initial condition,
\begin{equation}
\sup_{t\ge0}\mathbb E\,|X(t)|^2<\infty.
\end{equation}
In particular,
\begin{equation}
\limsup_{t\to\infty}\mathbb E\,|X(t)|^2\le \frac{c_0+c_\sigma\sigma^2}{c_1\,\underline a},
\end{equation}
where $\underline a=\min_i a_i$.
\end{theorem}

\section*{DOI}
Volume III v1 DOI: \href{https://doi.org/10.5281/zenodo.18844134}{10.5281/zenodo.18844134}.
This v2 supersedes v1 for formal mathematical use.

\end{document}
